\section{Integral principal multiples over rationally connected bases}
\label{sec:torsion}

We now bound a common principal multiple of the actual generalized
adjoint, including its denominators.  This will supply uniform
trivializations on rationally connected bases.

\begin{proposition}\label{prop:rc-torsion}
Fix an integer $1\leq d\leq 4$, a DCC set
$I\subset [0,1]\cap\Q$, and a positive integer $p$.
There exists a positive integer $\ell=\ell(d,I,p)$ with the following
property.  Let $Z$ be a normal projective complex variety of dimension $d$
with a rationally connected smooth projective resolution.  Suppose that
$(Z,B+M_Z)$ is generalized klt, that $B\geq 0$ has coefficients in $I$,
and that the b-nef rational b-divisor $\mathbf M$ has $p\mathbf M$
b-Cartier.  If
\[
 D=K_Z+B+M_Z\sim_{\Q}0,
\]
then $\ell D$ is an integral principal divisor.  Consequently, for every
integer $a\geq1$, the divisor $a\ell D$ is principal and
\[
 h^0\bigl(Z,\OO_Z(a\ell D)\bigr)=1.
\]
The same integer can be chosen for all dimensions in any fixed subset of
$\{1,2,3,4\}$.
\end{proposition}

\begin{proof}
There are three steps. Global ACC gives a finite coefficient set and
a uniform positive lower bound for generalized log discrepancies.
Boundedness up to isomorphism in codimension one then bounds torsion
in the divisor class group. Finally we clear the coefficients of the
actual adjoint before applying that torsion bound.

\smallskip
\noindent\emph{Finite coefficients and uniform discrepancies.}
We first construct the small modifications and the one-divisor extractions
that will allow us to use global ACC.  If an exceptional prime divisor $E$
over $Z$ is marked, assume that its generalized log discrepancy is
$a\in(0,1)$.  Choose a sufficiently high projective log resolution
$g:W\to Z$ on which the nef data descend and on which $E$ appears, when
marked.  Write
\[
 K_W+B_W+M_W=g^*D.
\]
The support of $B_W$ is SNC, all its coefficients are less than one, and
$B_W$ need not be effective.

We can also arrange that there is an effective exceptional Cartier divisor
$F$ on $W$ such that $-F$ is $g$-ample and
$\Supp(B_W)\cup\Supp(F)$ is SNC.  To see this, construct a resolution by
blowups with centres over the singular locus of the normal variety $Z$, and
then resolve the nef model, the marked valuation, and the boundary by further
blowups.
The latter can be taken with centres of codimension at least two on smooth
models.  All the resulting exceptional divisors therefore have images of
codimension at least two on $Z$.  At each blowup the negative of its
effective Cartier exceptional divisor is relatively ample.  Adding
sufficiently large multiples of the pullbacks of the divisors chosen at the
preceding stages gives an effective exceptional Cartier divisor whose
negative is ample for the composite morphism.  Additional blowups resolving
the union of the two supports preserve this construction.

Choose a positive rational number $\delta$ sufficiently small that all
coefficients of $B_W+\delta F$ are still less than one and, when $E$ is
marked, $\delta\mult_E F<a$.  For a sufficiently positive ample Cartier
divisor $A$ on $Z$, the divisor
\[
 M_W+g^*A-\delta F
\]
is ample.  Choose a general effective rational divisor
$\Lambda\sim_{\Q}M_W+g^*A-\delta F$ by dividing a general member of a
sufficiently divisible very ample multiple.  Bertini's theorem and the SNC
condition on the union of the supports show that
$(W,B_W+\delta F+\Lambda)$ is sub-klt.
There is a rational principal divisor $P$ on $W$ such that
\[
 \delta F+\Lambda=M_W+g^*A+P.
\]
Since $g$ is birational, $P=g^*(g_*P)$.  Pushing forward gives
\[
 K_Z+B+g_*\Lambda=D+A+g_*P,
\]
and pulling this equality back gives
\[
 g^*(K_Z+B+g_*\Lambda)
       =K_W+B_W+\delta F+\Lambda.
\]
Thus, setting $\Xi=B+g_*\Lambda$, the pair $(Z,\Xi)$ is an ordinary
effective rational klt pair.  If $E$ is marked, the general divisor
$\Lambda$ does not contain $E$, so its ordinary log discrepancy is
\[
 a(E,Z,\Xi)=a-\delta\mult_E F.
\]
This lies in $(0,1)$ by the choice of $\delta$ and the original bound $a<1$.

Take a log resolution of $(Z,\Xi)$ dominating $W$, so that it exhibits
$E$ when marked.  Apply the extraction corollary
\cite[Corollary~1.4.3 and its proof]{BCHM2010} with the prescribed set
$\{E\}$, or with the empty set when no valuation is marked.  The klt pair
$(Z,\Xi)$ satisfies the lc hypothesis, and $\Xi$ itself can be used as the
auxiliary klt boundary.  When a divisor is marked, its discrepancy is less
than one, so the extra condition for selected discrepancy-one valuations is
vacuous.  We obtain
a birational morphism
\[
 h:Z^+\longrightarrow Z
\]
with $Z^+$ $\Q$-factorial and with exactly the prescribed exceptional prime
divisors.  The construction in the cited proof is a log terminal model over
$Z$, hence $h$ is projective.  In the empty case it is a small
$\Q$-factorial modification.

Restore the original generalized data on $Z^+$ by setting
\[
 K_{Z^+}+B^++M_{Z^+}=h^*D,
\]
with the same nef b-divisor $\mathbf M$.  The resulting pair is generalized
klt and its adjoint is rationally linearly trivial.  Moreover, $B^+$ is
effective: it is the strict transform of $B$, together with $(1-a)E$ in
the extraction case.  If a marked divisor is already a prime on $Z$, the
small modification given by the empty-list construction suffices.

We next apply global ACC to these effective generalized pairs, whenever
their boundary coefficients belong to a fixed DCC set.  There are two
cases.  If the nef trace is not numerically trivial on a model determining
$\mathbf M$, then on a common higher model the nef part has the form
\[
 M_W=\frac1p(pM_W),
\]
where $pM_W$ is a nef Cartier divisor that is not numerically trivial.
The generalized global ACC theorem
\cite[Theorem~1.6]{BirkarZhang2016} applies: the dimension is fixed, the
pair is projective generalized lc, its total adjoint is numerically
trivial, the boundary coefficients belong to the chosen DCC set, and the
single nef weight is $1/p$.  We adjoin this fixed weight to that DCC set.
It follows that the possible boundary coefficients belong to a finite
set.

If instead $M_W\equiv0$, take a common higher model
$\pi:\widetilde W\to Z^+$ on which the nef data descend.  Since $Z^+$ is
$\Q$-factorial, the trace $M_{Z^+}$ is $\Q$-Cartier.  The divisor
\[
 M_{\widetilde W}-\pi^*M_{Z^+}
\]
is exceptional and numerically trivial over $Z^+$.  Applying the
negativity lemma to it and to its negative shows that it is zero.
Consequently $M_{Z^+}\equiv0$: every curve on $Z^+$ is dominated by a curve
on $\widetilde W$, and we use the projection formula.  Since the nef data
are pulled back from $Z^+$, the generalized discrepancies are the ordinary
discrepancies of $(Z^+,B^+)$.  This is an ordinary klt pair with
$K_{Z^+}+B^+\equiv0$, so ordinary global ACC
\cite[Theorem~1.5]{HMX2014} again gives a finite set of boundary
coefficients.

Applying this discussion to the small modification with no marked divisor
first shows that the coefficients of $B$ lie in a fixed finite rational
set $I_0\subset I$.  It also shows that the original generalized pairs
are uniformly generalized $\varepsilon$-lc for some
$\varepsilon=\varepsilon(d,I,p)>0$.  Indeed, otherwise we could choose a
sequence of such pairs and marked valuations whose generalized log
discrepancies satisfy
\[
 1>a_1>a_2>\cdots>0,\qquad a_i\longrightarrow0.
\]
Perform the extraction just constructed when the marked valuation is
exceptional, and otherwise take the small modification.  The boundary
coefficients of the resulting pairs belong to
\[
 I\ \cup\ \{1-a_i:i\geq1\},
\]
which is a DCC set because the added sequence is strictly increasing.
Global ACC would force the coefficients $1-a_i$ to belong to a finite
set, a contradiction.

\smallskip
\noindent\emph{Bounded topology and class-group torsion.}
We have obtained the uniform singularity bound without choosing any
bound for the class group. We now turn this geometric information into
a finite topological list.

Birkar's boundedness theorem
\cite[Theorem~1.7]{BirkarBoundedCY} now applies to the original varieties
$Z$.  Its hypotheses are
projectivity, fixed dimension, generalized $\varepsilon$-lc singularities,
effective boundary, a real-linearly trivial generalized adjoint, and
rational connectedness.  The singularity bound was just proved, the
adjoint is even rationally linearly trivial, and rational connectedness follows
from that of the given smooth projective resolution.  Thus, for every $Z$,
the theorem gives a projective variety $V_Z$, chosen from one bounded set,
that is isomorphic to $Z$ in codimension one.  Replacing each $V_Z$ by its
finite normalization preserves this property: normalization is an
isomorphism over the common normal open subset, and the complement still
has codimension at least two.  These normalizations still form a bounded set,
as in \cite[Section~9.5]{BirkarBoundedCY}.  Hence there is a projective
morphism $\mathcal V\to S$, with $S$ of finite type over $\C$, such that for
every $Z$ some fibre $V_s:=\mathcal V_s$, $s\in S(\C)$, is normal and
isomorphic to $Z$ in codimension one.

We explain explicitly how this boundedness controls torsion.  Put
$U=Z_{\mathrm{reg}}$.  The groups
\[
 H_1\bigl(U(\C),\Z\bigr)
\]
are finitely generated and take only finitely many isomorphism types.
First, this group is invariant under isomorphism in codimension one of
normal varieties.  The common open subset, intersected with the regular
loci, is obtained from each regular locus by removing a closed subset of
complex codimension at least two.  Removing such a subset from a smooth
complex variety does not change its fundamental group: choose a smooth
stratification of the removed subset and perturb paths and homotopies,
relative to their prescribed boundaries, to be transverse to its strata.
All these strata have real codimension at least four, so the perturbed
paths and two-dimensional homotopies avoid them.  The assertion for
$H_1$ follows by abelianization.

Second, we prove the assertion for the selected fibres $V_s$.  Replace $S$
by its reduction and partition it into finitely many locally closed strata
over which the family is flat.  Cover these strata by finitely many affine
opens.  On any one of these opens, still denoted $S$, let
\[
 \mathcal U=\operatorname{Sm}(\mathcal V/S)
\]
be the full relative smooth locus.  Flatness and finite presentation give
\[
 \mathcal U_s=\operatorname{Sm}(\mathcal V_s/\C)=(V_s)_{\mathrm{reg}}
\]
for each selected normal fibre, since regularity and smoothness agree over
$\C$.  Embed $\mathcal V$ in $\PP^N_{\C}\times S$.  Identify $\PP^N(\C)$
with the real algebraic set of Hermitian matrices $H$ satisfying $H^2=H$
and $\tr H=1$, by $[z]\mapsto zz^*/(z^*z)$, where $z^*$ is conjugate
transpose.  With real and imaginary coordinates on the affine $S$, the
entire map
\[
 \mathcal U(\C)\longrightarrow S(\C)
\]
is then a continuous semialgebraic map between affine semialgebraic spaces.
The semialgebraic triviality theorem of Delfs and Knebusch
\cite[Theorem~6.4]{DelfsKnebusch1982} gives a finite partition of $S(\C)$
into semialgebraic pieces over each of which the whole inverse image is a
product.  Thus the selected regular loci have only finitely many
homeomorphism types.  Each semialgebraic fibre has the homotopy type of a
finite simplicial complex by
\cite[Theorem~2.1 and Propositions~2.4--2.5]{DelfsKnebusch1982}.
The asserted finite list of finitely generated groups $H_1$ follows.

For each individual $Z$, the group $\Cl(Z)$ is also finitely generated.
Indeed, let $r:Y\to Z$ be a smooth projective rationally connected
resolution.  Weil divisor pushforward gives a surjection
\[
 \Pic(Y)\longrightarrow\Cl(Z):
\]
strict transforms give surjectivity on divisors, and principal divisors
push forward to principal divisors.  The Albanese map of $Y$ is constant
on every rational curve and therefore constant on $Y$, since rational
curves connect general points.  Consequently $\Pic^0(Y)=0$, and
N\'eron--Severi finiteness makes $\Pic(Y)$ finitely generated.

Normality and projectivity give
$\Gamma(U,\OO_U^*)=\C^*$, and restriction gives
$\Pic(U)=\Cl(Z)$ because $U$ is regular and contains every
codimension-one point.  The $n$th-power map on $\C^*$ is surjective.
For every $n\geq1$, the \'etale Kummer sequence
\cite[Tag~03PK]{StacksProject} gives the first isomorphism below.
Riemann existence for finite \'etale covers
\cite[Expos\'e~XII, Th\'eor\`eme~5.1 and Corollaire~5.2]{SGA1} gives the
second by classification of finite abelian torsors and abelianization:
\begin{equation}\label{eq:rc-kummer}
 \Cl(Z)[n]
 \simeq H^1_{\mathrm{\acute et}}(U,\mu_n)
 \simeq
 \Hom\bigl(H_1(U(\C),\Z),\mu_n(\C)\bigr).
\end{equation}
For a fixed $Z$, the orders of the groups on the left are bounded as $n$
varies, because $\Cl(Z)$ is finitely generated.  If $H_1(U(\C),\Z)$ had
a free summand of positive rank, the orders of the groups on the right
would grow without bound.  Hence $H_1(U(\C),\Z)$ is finite.  Only finitely
many such finite groups occur by the preceding paragraph, so there is a
uniform integer $T$ divisible by all their exponents.  Every torsion class
$c\in\Cl(Z)$ is killed by $T$: take $n$ equal to its order and apply
\eqref{eq:rc-kummer}, whose right-hand side has exponent dividing $T$.

\smallskip
\noindent\emph{An integral principal multiple.}
Finally, $pM_Z$ is an integral Weil divisor, being the pushforward of a
Cartier divisor on a model determining the nef data.  Choose a uniform
integer $q$ divisible by $p$ and by all denominators in the finite set
$I_0$.  Then $qD$ is an integral Weil divisor.  The relation
$D\sim_{\Q}0$ makes the class $[qD]\in\Cl(Z)$ torsion: if
$nD=\Div(v)$ for some $n\geq1$ and $v\in\C(Z)^*$, then
$n(qD)=\Div(v^q)$.  Hence
$TqD$ is principal.  Taking $\ell=Tq$ proves the proposition for dimension
$d$.  Taking a common multiple of the finitely many resulting integers
handles the stated range of dimensions.  All positive multiples of
$\ell D$ are principal, and projectivity and integrality of $Z$ give the
last assertion about sections.
\end{proof}

\begin{corollary}\label{cor:relative-rc-torsion}
Fix a finite rational coefficient set $\Phi\subset[0,1]\cap\Q$.
There is a positive integer $r=r(\Phi)$ such that the following holds.
Let $(X,B)$ be a projective klt fourfold pair with effective boundary
having coefficients in $\Phi$, and suppose that
$K_X+B\sim_{\Q}0$.  If there is a contraction $f:X\to Z$ with
$\dim Z>0$ and with a rationally connected smooth projective resolution
of $Z$, then $r(K_X+B)$ is an integral principal divisor. This includes
$B=0$, and every positive multiple is also principal.
\end{corollary}

\begin{proof}
Apply Theorem~\ref{thm:relative-denominators} to $f$ with the zero divisor
as the rational pullback class.  It gives uniformly bounded integers
$p_0\mid p$, a rational function $\psi\in\C(X)^*$, and an exact divisor
identity
\begin{equation}\label{eq:rc-exact-relative}
 K_X+B+\frac1{p_0}\Div(\psi)=f^*D_Z,
 \qquad D_Z=K_Z+B_Z+M_Z,
\end{equation}
where $(Z,B_Z+M_Z)$ is generalized klt, $B_Z$ is effective with
coefficients in a fixed rational DCC set, and $p\mathbf M$ is b-Cartier
with nef data.  The divisor $D_Z$ is rationally linearly trivial.  This
also follows directly from $f^*D_Z\sim_{\Q}0$: choose a positive
multiple making $D_Z$ Cartier and its pullback linearly trivial, and use
$f_*\OO_X=\OO_Z$ and the projection formula to trivialize that multiple
on $Z$.

Proposition~\ref{prop:rc-torsion} gives a uniform integer $\ell$ and
$v\in\C(Z)^*$ such that $\ell D_Z=\Div(v)$.  Choose $r$ divisible by both
$\ell$ and $p$, so that $p_0\mid r$.  Multiplying
\eqref{eq:rc-exact-relative} by $r$ gives
\[
 r(K_X+B)
 =\Div\!\left((v\circ f)^{r/\ell}\psi^{-r/p_0}\right).
\]
Both exponents are integers, so this is an equality to the divisor of an
actual rational function and includes the integrality of the adjoint
multiple.  All constants depend only on $\Phi$, since the possible
dimensions of $Z$ form a finite
set and the denominator theorem supplies uniform $p_0,p$ and a fixed
DCC set.
\end{proof}
